\documentclass{article}

\usepackage{amsmath}

\begin{document}

{\bf Worksheet \#1; date: 08/26/2018}

{\bf MATH 53 Multivariable Calculus}

\begin{enumerate}
\item {\em True / False?} Propositions can be thought of as a variable which may be true or false.

\item {\em True / False?} In propositional logic, a logical operator is a function that takes two input proposition and outputs a truth value.

\item {\em (Rosen 1.1.13a--f)} Let $p$ and $q$ be the propositions
\begin{align*}
p: & \text{You drive over $65$ per hour.} \\
q: & \text{You get a speeding ticket.}
\end{align*}
Write these propositions using $p$ and $q$ and logical connectives (including negations).
\begin{enumerate}
\item You do not drive over $65$ miles per hour.
\item You drive over $65$ miles per hour, but you do not get a speeding ticket.
\item You will get a speeding ticket if you drive over $65$ miles per hour.
\item If you do not drive over $65$ miles per hour, then you will not get a speeding ticket.
\item Driving over $65$ miles per hour is sufficient for getting a speeding ticket.
\end{enumerate}

\item Consider this sentence: I come back on campus only if I have teaching or meeting.
\begin{enumerate}
\item Rephrase it in the form ``if $p$ then $q$''
\item What is the converse, contrapositive and inverse of your answer to (a).
\end{enumerate}

\item Solve the system of equations:
\[
	\left\{\begin{array}{r c l}
	x^2 - 4x + 3 & \ne & 0 \\
	x^3 - 5x^2 + 6x & \ge & 0
	\end{array}\right.
\]

\item Without using a truth table, explain why the following proposition is always false.
\[
	(p \oplus q) \wedge (q \oplus r) \wedge (r \oplus p)
\]

\item {\em (Rosen 1.3.10c, 12c)} Show that
\[
	[p \wedge (p \to q)] \to q
\]
is tautology
\begin{enumerate}
\item with a truth table;
\item without a truth table.
\item {\em (Extra)} What does this proposition being tautology mean?
\end{enumerate}

\item {\em (Rosen 1.3.25)} Show that $(p \to r) \vee (q \to r)$ and $(p \wedge q) \to r$ are logically equivalent. How does this help us in writing proofs?

\item {\em (Rosen 1.3.33)} Show that $(p \to q) \to (r \to s)$ and $(p \to r) \to (q \to s)$ are not logically equivalent.
\end{enumerate}

\end{document}